\section{The converse with rational two-torsion}\label{pw:section}

We prove the pointwise converse for curves with $E(\Q)[2]\ne0$.
The finite graph construction supplies binary families with a unit coefficient
at every nonzero address.  Cyclotomic interpolation settles the even case.
For the odd case we first establish the missing lower bound for the moving
coefficient system, then keep one nonvanishing quadratic companion fixed
throughout the interpolation.

Fix such a curve $E/\Q$.  We use the prime weights $w(h)$, normalized values
$\mathcal L(h)$, and finite $2$-primary Shafarevich--Tate lengths $s(h)$
from the preceding sections.  The corank $c_2(E^{(h)})$ is a different
quantity: a corank-zero base has finite $s(h)$ before analytic
nonvanishing is known.  All new primes are odd primes of good reduction.
A filter fixes the real sign and the local unit squareclasses at a finite
support containing $2N$, and may exclude any further fixed finite set.

For a signed parameter $h$, write $A(h)$ for the symbol at the positive
index $|h|$ on its chosen filter.  We will use two bounds attached to a
unit coefficient:
\begin{align}
 A(h)=1\text{ in the even system}\quad&\Longrightarrow\quad
 L(E^{(h)},1)\ne0,\qquad v(\mathcal L(h))\le 2w(h)+C,
 \label{pw:even-unit-interface}\\
 A(h)=1\text{ in the odd system}\quad&\Longrightarrow\quad
 a(E^{(h)})=1,\qquad j(P_h)\le w(h)+C.
 \label{pw:odd-unit-interface}
\end{align}
Here $P_h$ is the genus Heegner sum with a fixed analytic-rank-zero
partner $h_*$, and $j(P_h)$ is its free-line index valuation.
The first bound is available from the even coefficient construction and
Theorem~\ref{thm:cyclotomic-bound}.  We use the second only after
Proposition~\ref{prop:rational-odd-bound} has supplied the arithmetic lower
bound required to normalize the odd system.
The contraction and isolation identities are those of
Theorem~\ref{thm:symbol-rules}; their finite graph consequences were proved
in Section~\ref{gr:section}.

Write $j_h(k)=j(P_{|h|}(h))$ for the ring-class index over $K=\Q(\sqrt{k})$.
Theorem~\ref{thm:heegner-ring-bound} states that, if
$k<0$ is an odd fundamental discriminant with $(k,2Nh)=1$, has at most
$A$ prime factors, is square at $2Nh$, and defines a nonexceptional
imaginary field, and the twists $h,hk$ have analytic
orders one and zero, then
\begin{equation}\label{pw:ring-lower-interface}
 2j_h(k)\ge4w(h)+s(h)+s(hk)-C_{E,A}.
\end{equation}
All coefficient and interpolation constants depend on the fixed curve,
its parametrization, and fixed local data.  A bound on $\omega(k)$ may
enter $C_{E,A}$; the sizes of its primes do not.  The finite networks,
required coefficient precision, and prime sizes may depend on the binary
dimension.  The unit bounds in Equations~\eqref{pw:even-unit-interface}
and~\eqref{pw:odd-unit-interface} may not.

By Lemma~\ref{pr:invariance}, an isogeny preserves analytic rank and
Selmer corank.  If the isogeny class contains a curve with full rational
two-torsion, we use that curve.  Otherwise choose a rational two-isogeny $E\to E'$ and write
$F=\mathbb Q(E[2])$, $F'=\mathbb Q(E'[2])$; both fields are quadratic.
A prime split in $F$ is called split, and one inert in $F$ is called simple.
Their weights are respectively $1$ and $1/2$.  A simple prime is good if it
is also inert in $F'$.  At such a prime
\begin{equation}\label{pw:good-simple}
 \#E(\mathbb F_p)\equiv2\pmod4.
\end{equation}
Indeed there is exactly one nonzero rational two-torsion point in the
reduction.  If the group order were divisible by four, its two-primary
subgroup would contain a rational point of order four.  Its image under the
two-isogeny would supply on $E'$ a rational two-torsion point different from
the dual kernel, contradicting inertness in $F'$.  Thus the good-twin rule
applies.  When $F=F'$, every simple prime is good.

\subsection{Seeds and finite supplies}\label{pw:seeds-section}

The word ``label'' below means an individually realizable fixed-support
elementary Frobenius label, with orientations when required.  It does not
include a self-residue of a prime generator.  All contractions and
comparisons are those of Theorem~\ref{thm:symbol-rules} and
Lemma~\ref{gr:forest}.

\begin{lemma}[Minimum seeds]\label{pw:minimum-seeds}
\begin{enumerate}
\item In an all-split alphabet, suppose a unit symbol has a nonempty
witness.  At the least positive vertex count $r$ of such a witness, all
mutual edge bits are immaterial.  The symbol is determined by the individual
fixed-support labels.  This remains true if the empty configuration is a
unit, provided no contraction to the empty configuration is used.
\item In the quadratic case, suppose the symbol has a nonzero witness,
the alphabet is closed under the allowed contractions, and it contains
the good twin operations.  There are
$r\ge0$ and $s_1\ge1$ such that for every $s\ge s_1$ of the permitted simple
parity, the minimum split count of a unit with $s$ simple vertices is $r$.
At these counts all edge bits are immaterial.  A chosen minimum seed may be
padded by arbitrarily many good twin pairs.
\end{enumerate}
The assertions also hold for a stabilized moving odd symbol, with every
particular finite list of uses made at a sufficiently advanced stage.
\end{lemma}

\begin{proof}
For the first assertion, a derivative in an edge of an $r$-vertex graph is a
symbol on $r-1$ vertices.  If $r\ge2$, this graph remains nonempty and has
value zero by minimality.  If $r=1$, there is no edge.  The omitted-vertex
description then leaves exactly the fixed-support labels.

For the second assertion let $r(s)$ be the minimum split count at simple
count $s$, with value $+\infty$ when no witness exists.  Once a finite value
occurs, good-twin insertion gives $r(s+2)\le r(s)$.  A nonincreasing sequence
of nonnegative integers stabilizes.  Choose $s_1\ge1$ in that range.  A
derivative of an edge with a split endpoint reduces the split count and
preserves the simple count; it is therefore zero at the minimum.  A
simple--simple derivative is zero by the symbol rule.  Twin insertion gives
the last assertion.

For the moving system, take the minimum and the nonincreasing sequence in
the stabilized profile.  Each comparison just used involves bounded vertex
counts.  A failure on an ultrafilter-large set of stages at any one such bound
would produce a bounded-weight sequence contradicting that minimum.  Thus
any fixed finite collection holds on an ultrafilter-large intersection of
stages, from which an arbitrarily high stage may be chosen.  No assertion
about every sufficiently high stage, or about all unbounded counts at one
stage, is needed.
\end{proof}

Recall that a list is partitionable when a permitted contraction forest
has a unit symbol as its output.  Its terminal polynomial is then nonzero,
since the indicated iterated finite difference is nonzero.

\begin{lemma}[Finite label reserves]\label{pw:reserves}
Fix a minimum seed and a finite label alphabet.
\begin{enumerate}
\item In the all-split rational alphabet, a list with the same total label
as the seed and sufficiently many copies of its labels is partitionable.
\item In the quadratic alphabet, retain $r$ split anchors with the seed
labels.  Suppose there are sufficiently many simple vertices of the required
parity and sufficiently many split reserves of each realizable rational
shift.  If $F\ne F'$, also suppose
\begin{equation}\label{pw:donor-balance}
 \begin{aligned}
 \Delta={}&\#\{p:\ p\text{ split in }F\text{ and inert in }F'\}\\
       &-\#\{p:\ p\text{ inert in }F\text{ and split in }F'\}>C_0.
 \end{aligned}
\end{equation}
Here $C_0$ depends only on the anchors and the reserves.  Then all other
split vertices can be attached to simple vertices so that the result is
the minimum seed and good twin pairs.  The twins may match a prescribed
finite enrichment of the rational labels.  The base simple labels may be
specified on any projection on which the required total agrees.
\end{enumerate}
All reserve requirements depend only on the finite alphabets, seed, and
projection.  They do not depend on the original list length.
\end{lemma}

\begin{proof}
In the first case retain the first $r-1$ seed labels and merge all other
vertices to the last cluster.  Sufficient copies ensure that this cluster
is nonempty.  Its label is forced by the total, hence is the last seed
label.  Lemma~\ref{pw:minimum-seeds} makes its unknown edges immaterial.

For the second case, the rational labels of split attachments form the
subgroup with trivial $F$-coordinate.  They act transitively on the simple
coset, so every prescribed change of a simple label has a realizable split
shift.  Reserve sufficiently many representatives of every such shift.
Remove these reserves and the anchors from the donor pool.  If $F\ne F'$,
\eqref{pw:donor-balance}, with $C_0$ exceeding the number removed, leaves a
distinct donor for each bad simple vertex.  Attach these donors; every
simple vertex is now good.  If $F=F'$, this operation was unnecessary.

Pair equal enriched good labels.  By withholding a bounded number of pairs,
leave between $s_1$ and $s_1+C_1$ unpaired vertices, with the correct parity;
$C_1$ depends only on the number of enriched labels.  Use reserved shifts
to turn the excess unpaired vertices into good pairs and to set the first
$s_1-1$ base positions to the desired seed labels.  The seed labels
themselves need not all be good.  Attach every remaining nonanchor split
vertex to the last base position.  Its projected label is forced by the
total, the anchors, the other base positions, and the twin pairs.  It is
the required last seed label.  Each cluster contains one simple vertex,
so every contraction is permitted.  Only a bounded number of reserve
shifts was used.

When incidences to a retained split anchor are needed only to remove twins,
they can be prescribed after the contracted labels are fixed, by the
independent terminal-edge realization.  Fixed stationary vertices instead
enter the finite label enrichment from the outset.  These two uses do not
increase the number of label types with the original list length.
\end{proof}

\begin{lemma}[Neutral packets]\label{pw:packets}
For any of the finite reserve problems above with prescribed rational total
zero, one can choose a finite packet of terminal labels with rational total
zero such that every nonempty union of copies of the packet has the required
supplies. In the quadratic case with $F\ne F'$, its donor surplus may
be made arbitrarily large. Compatible fixed splitting conditions may be
included whenever
the needed individual labels and shifts remain realizable.
\end{lemma}

\begin{proof}
Choose ample copies of every required label, including anchors and shifts.
When the quadratic fields $F,F'$ are distinct, add enough donors to
exceed all the losses from the other labels and the prescribed reserve
margin. When $F=F'$, every simple vertex is already good. Duplicate every
full label in the resulting list.  Duplication gives total zero in every
rational elementary coordinate and preserves any strict donor surplus.
Choose the margins so that after reserving the finitely many anchors in
any one packet, the union of any positive number of packets still meets
the reserve requirements.  All labels are realized by distinct fresh
primes; duplication never means repeating an actual prime factor.
\end{proof}

\subsection{The even converse and its uniform estimate}
\label{pw:lower-section}

\begin{proposition}[Even original base]\label{pw:even-base}
If $c_2(E)=0$, then $a(E)=0$.
\end{proposition}

\begin{proof}
Parity gives even functional sign.  Fix an even symbol on the filter of $1$,
including all auxiliary class-group tests.  If its value at $1$ is a
unit, analytic nonvanishing is already detected.  Otherwise use a
nonempty minimum seed in the full-torsion case, or a stabilized quadratic
seed in the other case.

For each $b$, assign a neutral packet from Lemma~\ref{pw:packets} to
every nonzero linear form on $\mathbb F_2^b$, and activate its primes
by that form.  At every nonzero address at least one packet is active;
their union satisfies the reserve problem and has a nonzero terminal
polynomial.  At zero no terminal is active.  The single-symbol case of
Theorem~\ref{gr:address} adds primary auxiliaries and makes the symbol a
unit at every nonzero address.  All local classes remain those of $1$.
Equation~\eqref{pw:even-unit-interface} gives the uniform bound at all
nonzero vertices.  Theorem~\ref{thm:cyclotomic-interpolation}, applied
at the original corank-zero base, now gives $a(E)=0$.
\end{proof}

\begin{proposition}[Partitionable even estimate]\label{pw:partitionable-even}
Fix an even unit symbol on a prescribed filter.  For every
partitionable parameter $h$ on its filter with $c_2(E^{(h)})=0$,
\begin{equation}\label{pw:partitionable-bound}
 \mathcal L(h)\ne0,\qquad v(\mathcal L(h))-2w(h)-s(h)\le C.
\end{equation}
The constant depends only on the symbol and fixed local data.
\end{proposition}

\begin{proof}
If the parameter has no prime factors, the only contraction is the identity,
so partitionability means that the empty symbol itself is a unit.
Equation~\eqref{pw:even-unit-interface}
then proves the assertion directly.  Otherwise the stationary terminal
set is nonempty.
Keep the prime factors of $h$ stationary.  The iterated edge difference
along the defining contraction forest is $1$, so their terminal polynomial
is nonzero.  The stationary one-symbol case
of Theorem~\ref{gr:address} gives arbitrarily large binary cubes based at
$h$, with only primary new primes, whose other addresses have unit symbol.
Equation~\eqref{pw:even-unit-interface} and $s(h_x)\ge0$ give the uniform
upper bound required by the cyclotomic transfer.  The base corank is zero,
so Theorem~\ref{thm:cyclotomic-interpolation} gives \eqref{pw:partitionable-bound}.  A base-dependent
dimension threshold does not affect the asserted constant.
\end{proof}

\subsection{The lower bound needed for odd coefficients}

The even estimate also supplies the arithmetic hypothesis of
Proposition~\ref{prop:odd-coefficients}. For an already analytic-rank-one
parameter $h$, we choose a companion $k'=k'(h)$ with boundedly many prime
factors and estimate the central value at $hk'$. Comparing the ring and
genus Heegner sums then bounds $j(P_h)$ from below. The companion in this
argument may vary with $h$; the odd coefficient system has not yet been
normalized.

To compare the two fixed isogenous curves $E,E'$, add subscripts to
$s,w,\mathcal L$. At a parameter $h$ already known to have analytic rank
zero, both twists have finite Tate--Shafarevich groups by the forward
theorem. Cassels's isogeny formula for the arithmetic BSD quotient
\citep{Cassels1965}, as explained in
\citet[Theorem~2.1 and the following discussion]{Tate1966}, gives
\begin{equation}\label{pw:isogeny-weight-comparison}
 s_E(h)+2w_E(h)=s_{E'}(h)+2w_{E'}(h)+O_{E,E'}(1).
\end{equation}
Indeed, at an active odd good prime the component-group valuation is
$t_p$ for the corresponding curve, so the active Tamagawa valuations
sum to twice its weight. At the fixed support only finitely many local
twisting classes occur. Differential, period, and torsion ratios are
bounded by the fixed isogeny and its reverse, and the rank-zero regulator
is $1$. These account for the bounded error in
\eqref{pw:isogeny-weight-comparison}. The isogeny identity applies to the
known finite groups and does not assume the BSD value formula.

The two twists have the same $L$-function. Their fixed normalizing periods
therefore give
$v(\mathcal L_E(h))-v(\mathcal L_{E'}(h))=O_{E,E'}(1)$, uniformly in $h$.
Combining this with \eqref{pw:isogeny-weight-comparison} shows that
$v(\mathcal L(h))-2w(h)-s(h)$ changes by a uniformly bounded amount.
Consequently \eqref{pw:partitionable-bound} transfers in either direction
with one constant, even when the choice between the two fixed isogeny
orientations depends on $h$.

\begin{lemma}[A bounded companion for a varying odd parameter]
\label{pw:rough-companion}
Fix a positive odd-sign filter and a negative genus partner $h_*$ with
$L(E^{(h_*)},1)\ne0$, chosen using Proposition~\ref{prop:rough-nonvanishing} and then
excluded from the variable parameters.  There are constants $A,C$ such that every $h$ on this filter of analytic
order one admits a negative odd fundamental discriminant $k'$ for which
\begin{enumerate}
\item $(k',2Nh)=1$, $k'$ is a square at every place of $2Nh$, and
$\mathbb Q(\sqrt{k'})$ is nonexceptional;
\item $\omega(k')\le A$;
\item $L(E^{(hk')},1)\ne0$ and
\begin{equation}\label{pw:rough-bound}
 v(\mathcal L(hk'))\le 2w(hk')+s(hk')+C.
\end{equation}
\end{enumerate}
\end{lemma}

\begin{proof}
Choose the class-group support and all other fixed label tests before
varying $h$.  Prepare an even negative-index symbol for each possible
total type at this support, and for both isogeny orientations if needed.
There are finitely many such preparations.  The positive index $|hk'|$
has nonzero total rational type: its negative-discriminant parity is
already nontrivial among the tests at $2$.  Thus in the full-torsion case
a unit seed has positive count.  In the quadratic case use the stabilized
seed of Lemma~\ref{pw:minimum-seeds}.

When $F\ne F'$, interchanging the isogeny orientation interchanges the two
counts in \eqref{pw:donor-balance}.  Select the orientation with initial
balance at least zero.  The existing donors can then pay for arbitrarily
many existing bad simples.  They do not have to be supplied by new primes.

Let $R$ be the absolute prime-count bound in Proposition~\ref{prop:rough-nonvanishing}.  Insert
in $k'$ a finite dummy supply that meets the anchors, minimum simple count,
and reserve requirements for every one of the finitely many prepared
symbols. When the quadratic fields $F,F'$ are distinct, add donor slack
paying for every reserved loss, the new dummy bad simples, and a further
possible loss of $R$ from the later rough factor. Add a fresh ramified
prime if needed to exclude the exceptional
imaginary fields.  The number of dummy primes is bounded independently
of $h$.  Their labels and initial incidences may be chosen by
Lemma~\ref{gr:realization}.

Absorb $h$ and the dummy product into the curve.  Apply
Proposition~\ref{prop:rough-nonvanishing} to complete $k'$ by a squarefree factor coprime to all
primes already used.  Prescribe its sign and its local unit classes so
that the final $k'$ is an odd negative fundamental discriminant, is a
square at $2Nh$, and has the required original local classes at $2N$.
These conditions are reduced congruences at the level of the absorbed
curve.  Additional fixed primes may be excluded.  The lemma makes the
central value at $hk'$ nonzero with at most $R$ new factors.  Its size
threshold may depend on $h$ and all these congruences; $R$ does not.

No particular resulting unit class at the auxiliary class-group support
was assumed in advance: dummy supplies were prepared for every possible
total there.  For the actual total, Lemma~\ref{pw:reserves} makes $hk'$
partitionable.  Analytic nonvanishing gives Selmer corank zero by the
forward theorem.  Proposition~\ref{pw:partitionable-even} therefore gives
\eqref{pw:rough-bound}, with the isogeny comparison above if the other
orientation was used.  The finite preparation and the bound on dummy and
rough factors give $A,C$ independent of $h$.
\end{proof}

Proposition~\ref{prop:heegner-comparison} compares the genus sum $P_h$
with the ring-class sum over $K=\Q(\sqrt{k})$:
\begin{equation}\label{pw:height-ratio}
 2j_h(k)-2j(P_h)=v(\mathcal L(hk))+O(1)
\end{equation}
whenever the two twists have analytic orders one and zero.
The error is uniform in $h$ for fixed $k$, and uniform in $h,k$ when
$\omega(k)$ is bounded and the prescribed local classes are retained.
Thus the same comparison applies to the bounded companions just constructed.

\begin{proposition}[Uniform lower bound for the odd detector]
\label{pw:odd-lower}\label{prop:rational-odd-bound}
Every analytically rank-one $h$ in the fixed positive filter satisfies
\begin{equation}\label{pw:odd-lower-bound}
 2j(P_h)\ge 2w(h)+s(h)-C.
\end{equation}
In particular $j(P_h)\ge w(h)-C$ uniformly in $h$ and its prime count.
\end{proposition}

\begin{proof}
Choose $k'$ by Lemma~\ref{pw:rough-companion}.  Combine
\eqref{pw:ring-lower-interface}, \eqref{pw:height-ratio}, and
\eqref{pw:rough-bound}.  Since the factors are coprime,
$w(hk')=w(h)+w(k')$, and $w(k')\le\omega(k')\le A$.  Hence
\begin{align*}
 2j(P_h)
 &\ge 4w(h)+s(h)+s(hk')-v(\mathcal L(hk'))-C\\
 &\ge 2w(h)+s(h)-2w(k')-C\\
 &\ge 2w(h)+s(h)-C'.
\end{align*}
Only even symbols and independently nonvanishing companions were used.
Thus this proof precedes, and does not assume, the odd minimum-depth
normalization.
\end{proof}

\subsection{A fixed companion for the odd converse}
\label{pw:fixed-companion-section}

Suppose the functional sign is odd. By
Proposition~\ref{pw:odd-lower}, the odd coefficient system can now be
normalized on the positive filter of $1$.  Fix its old support and genus
partner $h_*$.  Its moving detecting prime $\ell$ has not become an
extra fixed support requirement.

Choose a separate even symbol for negative discriminants square at $2N$,
excluding the fixed odd support, the relevant torsion ramification, and
the exceptional imaginary fields.  Choose a unit $k$ of that symbol.
Then $L(E^{(k)},1)\ne0$.  Fix this $k$ and
$K=\mathbb Q(\sqrt{k})$ for the rest of the construction, before choosing any
cube dimension.  Its fresh ramification makes it independent of the old
fixed label extensions.  At sufficiently advanced stages $\ell$ avoids
$k$.  The intrinsic odd label tests continue to omit the primes of $k$.
When the odd symbol already detects the base, the converse follows directly.
We henceforth consider the remaining case, in which the odd symbol has the
nonempty seed described above.

For the moving odd source, fix its stabilized minimum depth and seed
profiles before constructing packets. A bounded-weight seed has boundedly
many vertices here, since every prime has weight at least $1/2$. For each
fixed field used in the labels, permitting ramification above the one
moving prime $\ell$ adds only boundedly many dimensions to its elementary
squareclass space. Identify these spaces compatibly with their fixed
projections and stabilize the finitely many seed and label profiles.
The reserve bounds of Lemma~\ref{pw:reserves} are uniform over these
profiles, so the packet constructions below have a fixed size. They are
finite tables of labels; their prime representatives will be chosen at
an actual coefficient stage after the networks and their weight bounds
have been specified.

Our remaining task is to construct, for every binary dimension, a family
in which the odd test at $h$ and an even test at $hk$ are simultaneously
units away from the zero address. Every new prime will split in this
single field $K$. For quadratic $F$, the good-twin rule makes the unit
at $k$ a nonzero even test on additional prime configurations. With full
rational two-torsion, the restriction to $K$-split primes can instead
leave only the empty configuration at the original minimum depth. The
next construction supplies a nonempty witness in exactly that case;
its extra division changes the valuation estimate by a fixed constant.

\subsection{A nonempty witness with full rational two-torsion}
\label{pw:split-section}

\begin{lemma}[Recurrence congruences]\label{pw:recurrence}
Let $p$ be odd, $a\equiv p+1\pmod4$, and let
\[
 b_0=1,\quad b_1=a-\kappa,\quad
 b_{s+1}=ab_s-pb_{s-1}.
\]
If $\kappa=\pm1$, then $b_{s+1}-b_s$ is even and
\[
 \frac{b_{s+1}-b_s}{2}\equiv\frac{a-\kappa-1}{2}\pmod2
 \quad(s\ge0).
\]
If $\kappa=0$, put $b_{-1}=0$; then $b_s-b_{s-1}$ is odd for every
$s\ge0$.  Finally, for $\kappa=1$, if $a\equiv p+1\pmod{2^m}$ then
$b_s\equiv p^s\pmod{2^m}$ for every $s\ge0$.
\end{lemma}

\begin{proof}
Modulo $2$ the recurrence is $b_{s+1}=b_{s-1}$.  Its initial values prove
the parity assertions.  For $d_s=b_{s+1}-b_s$, the recurrence for $d_s$
is the same.  Directly,
$d_1-d_0\equiv(p-1)(p-\kappa)\equiv0\pmod4$ when $\kappa=\pm1$.
Thus $d_s/2$ is constant modulo $2$.  The last assertion follows by
substituting the sequence $p^s$ in the recurrence and checking its first
two terms.
\end{proof}

\begin{lemma}[Nonempty $K$-split witness]\label{pw:split-witness}
Suppose $E[2]$ is rational.  Choose a negative unit $k$ of a fixed even
symbol for $E$, prime to all previous fixed support and to the ramification
support of $E[8]$.  Put $K=\mathbb Q(\sqrt{k})$, chosen nonexceptional.
Let $S_0$ contain $2N$, the primes of $k$, and primes generating the class
group of $K$.  On the positive filter of $1$ square at $S_0$, there is an
even unit test on nonempty products of primes split in $K$ with these
properties:
\begin{enumerate}
\item it has a nonempty witness;
\item it satisfies the oriented split contraction and primary-isolation
rules over $K$, for contractions with nonempty output;
\item a unit at $D$ implies
\begin{equation}\label{pw:split-unit-bound}
 L(E^{(kD)},1)\ne0,\qquad v(\mathcal L(kD))\le 2\omega(D)+C_k.
\end{equation}
\end{enumerate}
The proof uses Lemmas~\ref{tc:trace-algebra} and~\ref{tc:inertia-identity}
to construct a further divided coefficient test.
\end{lemma}

\begin{proof}
Let $a_0(n)$ be the coefficients of the original negative-index form for
$E$.  Since $k$ attains its least normalized depth,
$v(a_0(|k|D)/a_0(|k|))\ge\omega(D)$ whenever $kD$ lies on its filter.
Use the ordinary even coefficient form for $E^{(k)}$, normalized by
$a(1)=1$, on the positive filter of $1$ square at $S_0$.
The two exact comparisons in Equation~\eqref{cs:even-ratio} give
\[
 2v(a(D))
 =v\!\left(\frac{\sqrt D\,L(E^{(kD)},1)}{L(E^{(k)},1)}\right)
 =2v\!\left(\frac{a_0(|k|D)}{a_0(|k|)}\right).
\]
Indeed, the period and the factor $\sqrt{|k|}$ cancel in the old ratio,
and the new period cancels in the new ratio.  Enlarging $S_0$ only
restricts the allowed $D$, so the original minimum still gives
\begin{equation}\label{pw:split-basic-depth}
 v(a(D))\ge\omega(D)
\end{equation}
at every allowed squarefree $D$.

If equality holds at a nonempty product of $K$-split primes, use the usual
even bit, projected from oriented $K$-labels to rational labels.  It has
all the required properties by Theorem~\ref{thm:symbol-rules}.  We
therefore assume for the rest of the proof that equality never occurs
there.  By the rational-phase assertion in Proposition~\ref{prop:even-coefficients},
normalization by $a(1)=1$ makes each nonzero squarefree coefficient a
rational number times a root of unity.  Its valuation is an integer,
not just an element of the possibly ramified coefficient field's
discrete value group.  Thus strict inequality implies
\begin{equation}\label{pw:split-extra-depth}
 v(a(D))\ge\omega(D)+1
 \quad\text{for every nonempty allowed all-$K$-split }D.
\end{equation}

Write $g=\sum a(n)q^n$ for the form.  Its characters satisfy
$\psi/\eta=\chi_{-4}$; since the scalar character $\psi$ is even,
$\eta$ is odd.  For a set $U$ of $K$-split
primes put
\[
 B_U=2^{-|U|}\prod_{p\in U}(U_p-\eta(p)V_p)g,
 \qquad
 \Theta_J=\sum_{m>0,\ (m,S_0)=1}\eta(m)m q^{Jm^2}.
\]
All expansions retain the prescribed support sieves.  The ordinary divided
form is
\[
 H_U=\frac12\left(B_U-\sum_{J\subseteq U}B_U[J]\Theta_J\right).
\]
Here $B_U[J]$ denotes its coefficient at the squarefree index $J$.

\smallskip\noindent\emph{The mod-$8$ consequence.}
We first prove
\begin{equation}\label{pw:split-mod8}
 a_p(E^{(k)})\equiv p+1\pmod8
 \quad\text{for every allowed good }p\text{ split in }K.
\end{equation}
Put $f=H_{\{p\}}\theta$ and apply its trace at a Frobenius lift $\phi_p$,
testing coefficient $p^2$.  The terms $B_{\{p\}}\theta$ and
$\Theta_{\{p\}}\theta$ have the same character, whose $p$-part is
primitive quadratic.  Thus they have no $p$-old constituent.  On their
conductor-one ramified principal-series constituents the identity is
\[
 U_p^2t(\phi_p)=U_p^3+d(\phi_p)U_p.
\]
The remaining term $\Theta_\varnothing\theta$ is unraised, and obeys the
ordinary good-prime coefficient identity.  Consequently
\[
 (t(\phi_p)f)[p^2]=f[p^3]+(d(\phi_p)f)[p].
\]
Replace the trace by a fresh split prime $p'$ with the same finite data,
including every support-unit test.  At each rational quadratic character
unramified outside $S_0$, a Frobenius lift at $p$ has the character value
of $p$.  Hence the matched $p'$ has the same value, and $pp'$ has value
$1$.  Equivalently, $p'$ has the same unit squareclasses as $p$ at all
places of $S_0$, so their product is square at all of them, including $2$.
Since $\phi_p\in G_K$, the new prime is also split in $K$.
Its coefficient on the left tests $pp'$; by
\eqref{pw:split-extra-depth} it is zero modulo the maximal ideal after
the two divisions.  To check the determinant term, let $d_1,d_0$ be the
determinant scalars on the raised and unraised character spaces.  Then
\[
 d(\phi_p)f
 =d_1f+\frac{B_{\{p\}}[1]}2(d_1-d_0)
                  \Theta_\varnothing\theta.
\]
The character difference $(d_1-d_0)/2$ is integral.  Also
$f[p]\equiv H_{\{p\}}[p]=0$ modulo the maximal ideal, since
$\Theta_{\{p\}}[p]=1$, whereas
$(\Theta_\varnothing\theta)[p]\equiv\Theta_\varnothing[p]=0$
because $\theta\equiv1$ and $p$ is nonsquare.  Thus
$(d(\phi_p)f)[p]=0$ after reduction.  This cancellation uses the two
character spaces together.

On the right, the indices $p$ and $p^3$ in $B_{\{p\}}$ have underlying
squarefree index $1$; they involve only $a(1),a(p^2),a(p^4)$.  The
$J=\varnothing$ unary term has no coefficient at either nonsquare
index.  Thus no other squarefree coefficient contributes there.
For $a=a_p(E^{(k)})$, the phase-free square coefficients are
$1,a-1,a(a-1)-p$.  Hence the remaining coefficient is, up to a unit,
\[
 H_{\{p\}}[p^3]
 =\frac{a^2-2a-p+1-p(a-2)}4
 =\frac{(a-1)(a-p-1)}4.
\]
This quotient is rational. Its vanishing in the residue field therefore
gives a positive integral $2$-adic valuation, so it is divisible by $2$.
Full rational two-torsion makes $a$ even, so
$a-1$ is odd; this proves \eqref{pw:split-mod8}.  No individual total
filter condition on $p$ was needed: its matched prime $p'$ makes $pp'$
lie on the filter, and the square-index terms on the right come from $1$.

\smallskip\noindent\emph{One further division by $2$.}
Define
\begin{equation}\label{pw:four-divided-form}
 \widetilde H_U=\frac14
 \left(B_U-\sum_{J\subseteq U}B_U[J]\Theta_J\right).
\end{equation}
We verify integrality at every coefficient.  Write a test index uniquely
as $Jrm^2$, where $J\subseteq U$ and $r$ is squarefree and disjoint from
$U$.  The underlying squarefree coefficient is at
$D=(U/J)r$.  Removing the unit phases, a prime of $U/J$ contributes an odd
factor $b_s-b_{s-1}$ with $\kappa=0$, and a prime of $J$ contributes
$b_{s+1}-b_s$ with $\kappa=(D/p)$.  Thus
\begin{equation}\label{pw:shell-valuation}
 v(B_U[Jrm^2])\ge\omega(r)
\end{equation}
by \eqref{pw:split-basic-depth}, with one additional unit of valuation
when $D$ is nonempty and all its primes split in $K$.

If $r\ne1$ and all primes of $D$ split in $K$, the last inequality is
at least $\omega(r)+1\ge2$.  If $D$ contains an inert prime, then $r$
contains at least two: the total rational $K$-character is trivial on
the filter, and all primes of $U$ split.  Then
\eqref{pw:shell-valuation} is again at least $2$.  No unary term has
this squarefree part.  This proves divisibility by $4$ for $r\ne1$.

Suppose $r=1$ and $J\ne U$.  Then $D=U/J$ is nonempty and all split,
so $B_U[J]$ is even.  This cancellation can be checked without dividing
by a leading multiplier, which may be zero.  Factor from both shell
coefficients the common phase and the scalar
$2^{-|U|+|J|}a(U/J)$, whose valuation is at least $1$ by
\eqref{pw:split-extra-depth}.  The remaining factors are the odd
differences at $U/J$, the integral divided differences at $J$, and
the ordinary square-index factors outside $U$.  By
Lemma~\ref{pw:recurrence}, their product is congruent modulo $2$ to
the product of their leading values, even when one of those values
is zero modulo $2$.  Since $m$ is odd, multiplication by $m$ preserves
this congruence.  The common scalar supplies the second factor $2$.
Retaining the phase $\eta(m)$ therefore gives
\[
 B_U[Jm^2]\equiv B_U[J]\eta(m)m\pmod4.
\]
This is exactly the cancellation with $B_U[J]\Theta_J$.

Finally let $J=U$, so $D=1$.  At a prime outside $U$ the recurrence gives
$b_s\equiv p^s\pmod4$.  At a prime of $U$,
\eqref{pw:split-mod8} gives $b_s\equiv p^s\pmod8$, and hence
\[
 \frac{b_{s+1}-b_s}{2}
 \equiv \frac{a_p-2}{2}\,p^s\pmod4.
\]
Multiplying these congruences, again retaining the phases, gives
$B_U[Um^2]\equiv B_U[U]\eta(m)m\pmod4$.  This proves integrality in
the last case and establishes \eqref{pw:four-divided-form}.
Diamonds preserve this lattice: when $J\ne U$, $B_U[J]$ is even and
the quadratic character difference supplies another factor $2$; when
$J=U$, the characters agree exactly.

\smallskip\noindent\emph{The new symbol and its witness.}
Multiply \eqref{pw:four-divided-form} by ordinary theta and use the
integral trace operators.  Two fresh traces in $G_K$ give zero after
reduction at every coefficient.  To see this, start with an arbitrary
test index $Jrm^2$ and choose the two fresh primes outside that index
and the level.  The filter imposes $\chi_K(D)=1$; since $U$ is all split,
the number of inert factors in $r$ is either zero or at least two.
The fresh traces replace $r$ by $rq_1q_2$, with both $q_i$ split.
In the first case the valuation after division by $4$ is at least
$\omega(r)+2+1-2=\omega(r)+1\ge1$.  In the second case it is at least
$\omega(r)+2-2=\omega(r)\ge2$.  The unary terms disappear under these
fresh tests.  Additional fresh $G_K$ traces only improve these estimates,
so vanishing holds on the entire reduced trace orbit.  The trace identity
therefore gives
\[
 t(gu^2)=t(g)\qquad(g,u\in G_K)
\]
on this reduced trace orbit.  Thus the tests factor through the elementary
extension of $K$ unramified off the level.

A single fresh $K$-split trace at coefficient $1$ reads the bit
\begin{equation}\label{pw:deeper-bit}
 \widetilde A(D)=
 \left(2^{-\omega(D)-1}a(D)\text{ with its phase removed}\right)\bmod2
\end{equation}
for a nonempty all-split configuration.  The contraction rule follows
from the same integral local identity as in Lemma~\ref{tc:inertia-identity}.
For completeness its first shell, with $p\in U$ and a fresh trace $q$,
has underlying configuration $D'=(U/p)q$, which is still nonempty, and
its reduction is
\[
 c_p(D')\widetilde A(D'),\qquad
 c_p(D')=\frac{a_p-(D'/p)-1}{2}\bmod2.
\]
Changing the oriented inertia bit changes $c_p$ once.  The omitted-prime
identity gives the product label and removes the entire rational prime
$p$ from the new ramification support.  Consequently its repeated
contraction has precisely the forest interpretation of
Lemma~\ref{gr:forest}.  An omitted-vertex label equal to the identity has
trace zero, which is primary isolation.  No contraction to the empty
configuration has occurred.

It remains to prove that this deeper test has a witness.  Fresh
ramification gives $K\not\subseteq\mathbb Q(E[8])$.  By Chebotarev there
is a good prime $p$ inert in $K$ and trivial on the original $E[8]$.
Then
\[
 p\equiv1\pmod8,\qquad a_p(E^{(k)})\equiv-2\pmod8.
\]
On $\widetilde H_\varnothing$, whose coefficient at $1$ is zero, test
$t(u)^2$ at coefficient $1$, where $u=\operatorname{Fr}_p$.  The square
coefficient and matching unary character give, up to a unit,
\[
 \frac{a_p(E^{(k)})-1-p}{4}\bmod2=1.
\]
The determinant term at coefficient $1$ cancels.  In characteristic two,
$t(u^2)=t(u)^2$, so a single trace in $G_K$ is nonzero.  A fresh prime
realizing this finite trace test is split in $K$.  Moreover $u^2$ is
trivial on every rational quadratic character, so the prime lies on the
required total rational filter.  It is a nonempty witness for
\eqref{pw:deeper-bit}.  Finally the exact coefficient ratio gives
\eqref{pw:split-unit-bound}, with the additional depth at most one.
\end{proof}

\begin{lemma}[Orientation correction]\label{pw:coinvariants}
Let $K$ be imaginary quadratic, let $S_0$ contain $2$ and its ramification
support, and let $V_0$ be the elementary label space over $K$ unramified
off $S_0$.  Complex conjugation acts on $V_0$ by $c$.  If the rational
total of an oriented prime configuration is zero on every rational
quadratic character unramified off $S_0$, then its oriented total belongs
to $(1+c)V_0$.  A reserved supply containing every label allows its
orientations to be flipped to realize any prescribed change in
$(1+c)V_0$.
\end{lemma}

\begin{proof}
Complex conjugation has order two outside $G_K$, so
$G_{\mathbb Q}=G_K\rtimes\langle c\rangle$.  A $c$-invariant quadratic
character of $G_K$ extends by giving $c$ value zero.  At a rational prime
outside $S_0$, inertia lies in $G_K$ and the character is trivial there;
thus its extension remains unramified outside $S_0$.  The assumed total
therefore annihilates $(V_0^*)^c$.

For a finite vector space with involution, the annihilator of
$\ker(1+c^*)=(V_0^*)^c$ is $\operatorname{im}(1+c)$.  This proves the
first assertion.  Flipping a vertex of label $v$ changes its oriented
label from $v$ to $cv$, and hence changes the total by $(1+c)v$.
Reserve one representative for every preimage needed to express the
desired element of the image.  Their independent flips give the second
assertion.  These flips preserve every rational norm label.
\end{proof}

\subsection{Common packets and rank-one interpolation}
\label{pw:assembly-section}

\begin{proposition}[Families with a fixed companion]
\label{pw:fixed-companion-families}
With the fixed odd coefficient system and companion $k$ chosen in
Section~\ref{pw:fixed-companion-section}, suppose the odd symbol at $1$
is zero. There is a constant $C$ with the following property. For every $b\ge1$ there are distinct primes outside
$2Nk$ and the fixed supports of both coefficient tests, with nonzero
linear activations, such that their signed products $h_x$,
$x\in\F_2^b$, satisfy $h_0=1$. Every $h_x$ is positive and has the fixed
local classes of $1$, and every new prime splits in $K=\Q(\sqrt{k})$.
The primes also avoid the detecting prime of the actual odd coefficient
stage. In particular $(h_x,k)=1$ for every $x$. For every $x\ne0$,
\[
 a(E^{(h_x)})=1,\qquad L(E^{(h_xk)},1)\ne0,
\]
and
\[
 j(P_{h_x})\le w(h_x)+C,\qquad
 v(\mathcal L(h_xk))\le 2w(h_x)+C.
\]
The constant $C$ and the discriminant $k$ are independent of $b$.
\end{proposition}

\begin{proof}
We construct common terminal packets, individually split in $K$, such
that the odd terminal polynomial and an even terminal polynomial testing
$hk$ are separately nonzero at each nonzero address.

First suppose $F,F'$ are quadratic.  Keep the original even symbol
evaluated on $hk$, treating the vertices of $k$ as stationary external
data.  Enrich the rational labels of the $h$-vertices by the primes of
$k$.  Their total there and at the even support is required to be zero.
Let $V$ be the resulting rational elementary label space, including the
$F,F'$ and $K$ coordinates and all relations among these characters.
The allowed simple labels form the affine set
$\{v\in V:\chi_F(v)=1,\ \chi_K(v)=0\}$, in additive notation.
The difference of any two of its elements belongs to
$\ker\chi_F\cap\ker\chi_K$, precisely the realizable rational labels
of the permitted split attachments.  Translation by this subgroup is
therefore transitive.  The fields $F,F'$ are unramified outside the old
support, whereas $K$ has a fresh ramified prime; thus $K$ is independent
of their compositum.  When $F\ne F'$, this realizes both the good-simple
pattern $(1,1,0)$ and the donor pattern $(0,1,0)$ in their three
coordinates.  When $F=F'$, all simple labels are good.

For the old odd projection, take a normal closure of its finite label
extension, including the moving prime $\ell$ at any stage realizing the
stabilized profile. It is unramified outside the old support and $\ell$.  Once $\ell$ avoids
$k$, a prime ramified in $K$ remains outside that set.  Hence $K$
has trivial intersection with this normal closure, and every old odd
label can be realized together with splitting in $K$.  This applies
also to oriented split anchors over $F$.  The added stationary $k$
coordinates may have relations, but they are already included in $V$;
they do not impose extra requirements on the omitted odd projection.

Use packets with the required simple supply and reserves in this kernel,
and with sufficiently large donor slack when $F\ne F'$. For the even
test, take no split anchors.
Correct all bad simples to good ones, pair equal enriched labels, and
retain a nonempty good base pair.  Use reserves to match the pair and
the remaining good pairs, and attach every leftover split vertex to the
last base position.  Its label is forced by total zero.  The output is
a collection of good twin pairs, each matching also at every split
vertex of $k$.  Twin insertion into the unit at $k$ proves that this
output has even symbol one.  Thus the even terminal polynomial is
nonzero.

For the odd test, project away the newly added tests at $k$.  The same
packets contain enough split anchors, simple vertices and reserve shifts
to contract to the stabilized odd seed and good pairs, by
Lemma~\ref{pw:reserves}.  Its total on that projection is the prescribed
zero total.  Incidences between the retained anchors and the good pairs
can be set after contraction, by independent edge choices.  The odd
polynomial is therefore nonzero as well.  This odd test does not include
the stationary vertices of $k$.

Now suppose the curve has full rational two-torsion.  Use the even test
of Lemma~\ref{pw:split-witness}, with $K$-oriented labels on its support
$S_0$, and let $v_1,\ldots,v_r$ be a minimum nonempty witness.  Its
oriented total $v_1+\cdots+v_r$ lies in $(1+c)V_0$, by
Lemma~\ref{pw:coinvariants}.  Choose packets containing ample copies of
every needed even label, all odd seed labels on their rational
projection, and orientation-shifting reserves.  In the common compositum,
restriction to the even label extension is surjective, so each even label
has a lift.  Each intrinsic odd label also has a lift fixing $K$: fresh
ramification gives $K\cap L_{\rm odd}=\mathbb Q$, and the moving odd prime
avoids $k$.  Extending these elements to the common compositum gives the
required lifts.  Take sufficiently many lifts for each projection
separately, without prescribing their other projection.  Duplicate each
full label, so the packet has rational total zero.

For the even contraction, retain anchors with labels
$v_1,\ldots,v_{r-1}$.  After reserving these anchors, use orientation
flips to make the total equal to $\sum v_i$; this is possible by
Lemma~\ref{pw:coinvariants}.  Merge all remaining vertices into the
last cluster.  Its label is $v_r$.  The cluster is nonempty, and the
minimum nonempty count makes the output a unit independently of its
mutual edges.  For the odd contraction, use its rational projection.
Every odd seed label lifts to a $K$-split prime because $K$ was freshly
ramified outside its intrinsic support.  Retain all but the last odd
seed label and merge the remainder.  Total zero gives the last label,
so this is also a unit.  Orientation correction is conjugation over
$\mathbb Q$, which preserves every rational quadratic label, including
those at the moving odd prime.  Adding $k$ to the common support has not
added it to the intrinsic odd projection.

In both cases put a sufficient common packet at every nonzero linear
form on $\mathbb F_2^b$.  Each nonzero address has two separately
nonzero terminal polynomials.  Theorem~\ref{gr:address} supplies primary
auxiliaries split in $K$ (and $F$ in the quadratic case), neutral at the
stationary primes of $k$, and makes both symbols one at every nonzero
address.  The auxiliary generators are primary at every support
coordinate used by either projection.  Thus the odd and even graph
identities remain identities on the same common terminal polynomial
data.  All active products $h$ are positive and square at the fixed
tests, and all their primes split in $K$.

For each fixed $b$, the terminal tables and networks now give a finite
list of coefficient tests and a maximum weight. Choose the working
precision after these bounds are fixed. Finer ray neighborhoods affect
prime realization but do not enlarge the elementary label alphabet.
Lemma~\ref{cs:fixed-weight-uniformity} and
Theorem~\ref{thm:symbol-rules} realize the required divisibilities and
symbol identities on an ultrafilter-large intersection of stages.
At an arbitrarily high actual stage in that intersection, realize the
finite graphs by distinct primes. The unit-depth constants and the
companion $k$ were fixed before $b$ was chosen, and therefore remain
uniform in these realizations.

At each nonzero address the odd unit detects analytic rank one and gives
$j(P_h)\le w(h)+C$. The even unit detects $L(E^{(hk)},1)\ne0$ and gives
$v(\mathcal L(hk))\le2w(hk)+C$, or the bound
\eqref{pw:split-unit-bound} in the full-torsion case. The fixed quantity
$2w(k)$ is absorbed into $C$. These are the asserted estimates.
\end{proof}

\begin{proposition}[Odd original base]\label{pw:odd-base}
If $c_2(E)=1$, then $a(E)=1$.
\end{proposition}

\begin{proof}
Parity gives odd functional sign. Proposition~\ref{pw:odd-lower}
therefore permits the odd normalization used above. If its symbol at $1$
is a unit, it already gives $a(E)=1$. Otherwise apply
Proposition~\ref{pw:fixed-companion-families} for arbitrarily large $b$.
At every nonzero address its estimates and
Proposition~\ref{prop:heegner-comparison}, with the fixed $k$, give
\[
 2j_h(k)\le4w(h)+C.
\]
The finite Sha lengths there are nonnegative. Hence
$2j_h(k)-4w(h)-s(h)-s(hk)\le C$, as required for Heegner interpolation.

At address zero, $L(E^{(k)},1)\ne0$ gives $c_2(E^{(k)})=0$ by
Lemma~\ref{pr:ranks}. Lemma~\ref{pr:invariance} gives
\[
 \corank_{\Ztwo}\Sel_{2^\infty}(E/K)
 =c_2(E)+c_2(E^{(k)})=1.
\]
Since $K$ is fixed, the finite quotient of this base Selmer group has a
fixed length. Theorem~\ref{thm:heegner-interpolation} makes
$L(E,s)L(E^{(k)},s)$ have a simple zero. Its second factor is nonzero at
$1$, so $a(E)=1$.
\end{proof}

\begin{theorem}[The rational-two-torsion converse]\label{thm:rational-converse}
Let $E/\Q$ satisfy $E(\Q)[2]\ne0$ and
$c_2(E)=\operatorname{corank}_{\Z_2}\Sel_{2^\infty}(E/\Q)\in\{0,1\}$,
with the local Kummer conditions at every place.  Then
\[
                 \operatorname{ord}_{s=1}L(E,s)=c_2(E).
\]
There is no restriction on the rational isogeny configuration, complex
multiplication, or reduction at $2$ or at the bad primes.
\end{theorem}

\begin{proof}
The initial isogeny choice preserves both quantities and covers both
rational two-torsion configurations.  Proposition~\ref{pw:even-base}
proves the corank-zero case.  In the corank-one case,
Proposition~\ref{prop:rational-odd-bound} supplies the coefficient
normalization, and Proposition~\ref{pw:odd-base} proves analytic rank one
using a companion fixed before the binary dimension grows.
\end{proof}
